\sectionheading{5.44 Python-এ Function}

একই ধরনের কাজ Program-এর বিভিন্ন স্থানে বারবার প্রয়োজন হলে একই Code
বারবার না লিখে সেই কাজকে একটি নির্দিষ্ট অংশে রাখা যায়।
এই পুনঃব্যবহারযোগ্য Code Block-কে \textbf{Function} বলা হয়।

%================================================
% Definition
%================================================

\booksubheading{Function কী?}

\begin{definitionbox}

নির্দিষ্ট কোনো কাজ সম্পন্ন করার জন্য তৈরি পুনঃব্যবহারযোগ্য
Code Block-কে \textbf{Function} বলে।

\end{definitionbox}

\begin{center}

\textbf{একবার Function তৈরি}
\(\longrightarrow\)
\textbf{প্রয়োজনমতো Call}
\(\longrightarrow\)
\textbf{একই কাজ পুনরায় সম্পন্ন}

\end{center}

%================================================
% Benefits
%================================================

\booksubheading{Function ব্যবহারের সুবিধা}

\begin{compacttwocolpoints}

\item একই Code বারবার লিখতে হয় না।

\item Code পুনঃব্যবহার করা যায়।

\item Program ছোট, গোছানো ও সহজবোধ্য হয়।

\item বড় Problem-কে ছোট ছোট অংশে ভাগ করা যায়।

\item Program পরিবর্তন, পরীক্ষা ও রক্ষণাবেক্ষণ সহজ হয়।

\end{compacttwocolpoints}

%================================================
% Types
%================================================

\booksubheading{Python-এ Function-এর ধরন}

\begin{tabularx}
{\linewidth}
{
|>{\bfseries\RaggedRight\arraybackslash}p{0.27\linewidth}
|Y
|>{\RaggedRight\arraybackslash}p{0.32\linewidth}|
}

\hline
\rowcolor{softBlue}

\textbf{ধরন}
&
\textbf{অর্থ}
&
\textbf{উদাহরণ}
\\
\hline

Built-in Function
&
Python-এ আগে থেকেই তৈরি।
&
\textenglish{\ttfamily print(), input(), len(), type()}
\\
\hline

User-defined Function
&
Programmer প্রয়োজন অনুযায়ী তৈরি করে।
&
\textenglish{\ttfamily add(), area(), greet()}
\\
\hline

\end{tabularx}

\begin{rememberbox}

\textenglish{\ttfamily print()},
\textenglish{\ttfamily input()},
\textenglish{\ttfamily len()},
\textenglish{\ttfamily type()}
ইত্যাদি হলো \textbf{Built-in Function}।

\textenglish{\ttfamily def}
ব্যবহার করে Programmer যে Function তৈরি করে তা
\textbf{User-defined Function}।

\end{rememberbox}

%================================================
% Syntax
%================================================

\booksubheading{User-defined Function তৈরির Syntax}

\begin{codeoutputbox}{Function-এর সাধারণ Structure}

\noindent
\begin{minipage}[t]{0.52\linewidth}
\codeheading

\begin{english}
\ttfamily
def function\_name():\\
\hspace*{4mm}statement
\end{english}

\end{minipage}
\hfill
\begin{minipage}[t]{0.44\linewidth}
\outputheading

\textenglish{\ttfamily def}
\(\longrightarrow\)
Function তৈরির Keyword

\medskip

\textenglish{\ttfamily function\_name}
\(\longrightarrow\)
Function-এর নাম

\medskip

\textenglish{\ttfamily ()}
\(\longrightarrow\)
Parameter-এর স্থান

\medskip

\textenglish{\ttfamily :}
\(\longrightarrow\)
Header-এর শেষ

\end{minipage}

\end{codeoutputbox}

Function Body অবশ্যই সঠিকভাবে \textbf{Indented} হতে হবে।

%================================================
% Define and Call
%================================================

\booksubheading{Function তৈরি ও Call}

\begin{codeoutputbox}{প্রথম Function}

\noindent
\begin{minipage}[t]{0.53\linewidth}
\codeheading

\begin{english}
\ttfamily
def greet():\\
\hspace*{4mm}print("Hello")\\[0.5mm]
greet()
\end{english}

\end{minipage}
\hfill
\begin{minipage}[t]{0.43\linewidth}
\outputheading

\textbf{Output:}

\begin{center}
\textenglish{\ttfamily Hello}
\end{center}

\textenglish{\ttfamily def greet():}
Function Define করেছে।

\textenglish{\ttfamily greet()}
Function-টিকে Call করেছে।

\end{minipage}

\end{codeoutputbox}

\begin{rememberbox}

শুধু Function Define করলে তার Body Execute হয় না।

\begin{center}

\textbf{Function Call}
\(\longrightarrow\)
\textbf{Function Body Execute}

\end{center}

\end{rememberbox}

%================================================
% Multiple Call
%================================================

\booksubheading{একই Function একাধিকবার Call}

\begin{codeoutputbox}{Function Reuse}

\noindent
\begin{minipage}[t]{0.53\linewidth}
\codeheading

\begin{english}
\ttfamily
def greet():\\
\hspace*{4mm}print("Welcome")\\[0.5mm]
greet()\\
greet()\\
greet()
\end{english}

\end{minipage}
\hfill
\begin{minipage}[t]{0.43\linewidth}
\outputheading

\begin{english}
\ttfamily
Welcome\\
Welcome\\
Welcome
\end{english}

\medskip

Function একবার Define করা হয়েছে,
কিন্তু তিনবার Call করা হয়েছে।

এটাই Code Reuse-এর একটি সুবিধা।

\end{minipage}

\end{codeoutputbox}

%================================================
% Parameter Argument
%================================================

\booksubheading{Parameter ও Argument}

Function-এ বাইরে থেকে Data পাঠানোর জন্য Parameter ব্যবহার করা হয়।

\begin{codeoutputbox}{Parameter ও Argument}

\noindent
\begin{minipage}[t]{0.53\linewidth}
\codeheading

\begin{english}
\ttfamily
def greet(name):\\
\hspace*{4mm}print("Hello", name)\\[0.5mm]
greet("Rafi")
\end{english}

\end{minipage}
\hfill
\begin{minipage}[t]{0.43\linewidth}
\outputheading

\begin{center}
\textenglish{\ttfamily Hello Rafi}
\end{center}

\textenglish{\ttfamily name}
\(\longrightarrow\)
Parameter

\medskip

\textenglish{\ttfamily "Rafi"}
\(\longrightarrow\)
Argument

\end{minipage}

\end{codeoutputbox}

\begin{tabularx}
{\linewidth}
{
|>{\bfseries\RaggedRight\arraybackslash}p{0.27\linewidth}
|Y|
}

\hline
\rowcolor{softBlue}

\textbf{বিষয়}
&
\textbf{অর্থ}
\\
\hline

Parameter
&
Function Definition-এর মধ্যে ব্যবহৃত Variable।
\\
\hline

Argument
&
Function Call-এর সময় Parameter-এর জন্য পাঠানো Actual Value।
\\
\hline

\end{tabularx}

\begin{rememberbox}

\begin{center}

\textbf{Definition-এর মধ্যে}
\(\longrightarrow\)
\textbf{Parameter}

\qquad

\textbf{Call-এর মধ্যে}
\(\longrightarrow\)
\textbf{Argument}

\end{center}

\end{rememberbox}

%================================================
% Multiple Parameters
%================================================

\booksubheading{একাধিক Parameter}

\begin{codeoutputbox}{দুইটি Parameter}

\noindent
\begin{minipage}[t]{0.53\linewidth}
\codeheading

\begin{english}
\ttfamily
def add(a, b):\\
\hspace*{4mm}print(a + b)\\[0.5mm]
add(10, 20)
\end{english}

\end{minipage}
\hfill
\begin{minipage}[t]{0.43\linewidth}
\outputheading

\begin{center}
\textenglish{\ttfamily 30}
\end{center}

\textenglish{\ttfamily a, b}
হলো Parameter।

\textenglish{\ttfamily 10, 20}
হলো Argument।

\end{minipage}

\end{codeoutputbox}

%================================================
% Positional
%================================================

\booksubheading{Positional Argument}

Argument-এর Position অনুযায়ী Parameter-এ Value গেলে তাকে
\textbf{Positional Argument} বলা হয়।

\begin{codeoutputbox}{Positional Argument}

\noindent
\begin{minipage}[t]{0.55\linewidth}
\codeheading

\begin{english}
\ttfamily
def student(name, age):\\
\hspace*{4mm}print(name, age)\\[0.5mm]
student("Rafi", 18)
\end{english}

\end{minipage}
\hfill
\begin{minipage}[t]{0.41\linewidth}
\outputheading

\begin{center}
\textenglish{\ttfamily Rafi 18}
\end{center}

প্রথম Argument
\textenglish{\ttfamily "Rafi"}
যায়
\textenglish{\ttfamily name}-এ।

দ্বিতীয় Argument
\textenglish{\ttfamily 18}
যায়
\textenglish{\ttfamily age}-এ।

\end{minipage}

\end{codeoutputbox}

\begin{rememberbox}

Positional Argument-এ Argument-এর
\textbf{Order গুরুত্বপূর্ণ}।

\end{rememberbox}

%================================================
% Keyword argument
%================================================

\booksubheading{Keyword Argument}

Parameter-এর নাম উল্লেখ করে Argument পাঠানো হলে তাকে
\textbf{Keyword Argument} বলা হয়।

\begin{codeoutputbox}{Keyword Argument}

\noindent
\begin{minipage}[t]{0.56\linewidth}
\codeheading

\begin{english}
\ttfamily
def student(name, age):\\
\hspace*{4mm}print(name, age)\\[0.5mm]
student(age=18, name="Rafi")
\end{english}

\end{minipage}
\hfill
\begin{minipage}[t]{0.40\linewidth}
\outputheading

\begin{center}
\textenglish{\ttfamily Rafi 18}
\end{center}

Parameter-এর নাম উল্লেখ করা হয়েছে।

তাই এখানে Argument-এর Position বদলালেও
সঠিক Parameter-এ Value গেছে।

\end{minipage}

\end{codeoutputbox}

%================================================
% Default parameter
%================================================

\booksubheading{Default Parameter}

Function Definition-এর সময় কোনো Parameter-এর Default Value
নির্ধারণ করা যায়।

\begin{codeoutputbox}{Default Value}

\noindent
\begin{minipage}[t]{0.54\linewidth}
\codeheading

\begin{english}
\ttfamily
def greet(name="Student"):\\
\hspace*{4mm}print("Hello", name)\\[0.5mm]
greet()\\
greet("Rafi")
\end{english}

\end{minipage}
\hfill
\begin{minipage}[t]{0.42\linewidth}
\outputheading

\begin{english}
\ttfamily
Hello Student\\
Hello Rafi
\end{english}

\medskip

প্রথম Call-এ Argument নেই,
তাই Default Value ব্যবহৃত হয়েছে।

দ্বিতীয় Call-এ
\textenglish{\ttfamily "Rafi"}
Default Value-এর জায়গায় এসেছে।

\end{minipage}

\end{codeoutputbox}

%================================================
% Return
%================================================

\booksubheading{\textenglish{\ttfamily return} Statement}

Function-এর Result Function-এর বাইরে ফেরত পাঠানোর জন্য
\textenglish{\ttfamily return}
ব্যবহার করা হয়।

\begin{codeoutputbox}{\textenglish{\ttfamily return} ব্যবহার}

\noindent
\begin{minipage}[t]{0.54\linewidth}
\codeheading

\begin{english}
\ttfamily
def add(a, b):\\
\hspace*{4mm}return a + b\\[0.5mm]
result = add(10, 20)\\
print(result)
\end{english}

\end{minipage}
\hfill
\begin{minipage}[t]{0.42\linewidth}
\outputheading

\begin{center}
\textenglish{\ttfamily 30}
\end{center}

\textenglish{\ttfamily return a + b}
Result Function-এর বাইরে পাঠিয়েছে।

ফেরত পাওয়া Value
\textenglish{\ttfamily result}
Variable-এ রাখা হয়েছে।

\end{minipage}

\end{codeoutputbox}

%================================================
% print vs return
%================================================

\booksubheading{\textenglish{\ttfamily print()} বনাম \textenglish{\ttfamily return}}

\begin{tabularx}
{\linewidth}
{
|>{\bfseries\RaggedRight\arraybackslash}p{0.27\linewidth}
|Y
|Y|
}

\hline
\rowcolor{softBlue}

\textbf{বিষয়}
&
\textenglish{\ttfamily print()}
&
\textenglish{\ttfamily return}
\\
\hline

প্রধান কাজ
&
Output Screen-এ প্রদর্শন করে।
&
Value Function-এর বাইরে ফেরত দেয়।
\\
\hline

Result সংরক্ষণ
&
শুধু Print করলে Return Value হিসেবে পাওয়া যায় না।
&
Variable-এ রাখা যায়।
\\
\hline

পরে Calculation
&
Printed Output সরাসরি Result হিসেবে ব্যবহার করা যায় না।
&
Return করা Value ব্যবহার করা যায়।
\\
\hline

Function শেষ
&
অবশ্যই Function শেষ করে না।
&
Execute হলে Function শেষ হয়।
\\
\hline

\end{tabularx}

\begin{codeoutputbox}{Return করা Value পুনরায় ব্যবহার}

\noindent
\begin{minipage}[t]{0.53\linewidth}
\codeheading

\begin{english}
\ttfamily
def square(n):\\
\hspace*{4mm}return n * n\\[0.5mm]
x = square(5)\\
y = x + 10\\
print(y)
\end{english}

\end{minipage}
\hfill
\begin{minipage}[t]{0.43\linewidth}
\outputheading

\begin{center}
\textenglish{\ttfamily 35}
\end{center}

\textenglish{\ttfamily square(5)}
ফেরত দেয় 25।

তারপর—

\[
25 + 10 = 35
\]

\end{minipage}

\end{codeoutputbox}

\begin{rememberbox}

\textbf{\color{deepOrange}MCQ Trap:}

\begin{center}

\textenglish{\ttfamily print()}
\(\neq\)
\textenglish{\ttfamily return}

\end{center}

\textenglish{\ttfamily print()}
Output দেখায়, আর
\textenglish{\ttfamily return}
Result ফেরত দেয়।

\end{rememberbox}

%================================================
% Return stops function
%================================================

\booksubheading{\textenglish{\ttfamily return} হলে Function কোথায় থামে?}

\textenglish{\ttfamily return}
Execute হলে Function-এর Execution সেখানেই শেষ হয়।

\begin{codeoutputbox}{\textenglish{\ttfamily return}-এর পরের Code}

\noindent
\begin{minipage}[t]{0.52\linewidth}
\codeheading

\begin{english}
\ttfamily
def test():\\
\hspace*{4mm}print("A")\\
\hspace*{4mm}return\\
\hspace*{4mm}print("B")\\[0.5mm]
test()
\end{english}

\end{minipage}
\hfill
\begin{minipage}[t]{0.44\linewidth}
\outputheading

\begin{center}
\textenglish{\ttfamily A}
\end{center}

\textenglish{\ttfamily return}
Execute হওয়ার পর Function শেষ হয়েছে।

তাই
\textenglish{\ttfamily print("B")}
Execute হয়নি।

\end{minipage}

\end{codeoutputbox}

\begin{rememberbox}

\begin{center}

\textenglish{\ttfamily return}
\(\longrightarrow\)
\textbf{Value ফেরত দিতে পারে}
+
\textbf{Function Execution শেষ করে}

\end{center}

\end{rememberbox}

%================================================
% None
%================================================

\booksubheading{Explicit \textenglish{\ttfamily return} না থাকলে কী হয়?}

Python Function-এ Explicit
\textenglish{\ttfamily return}
না থাকলে Function সাধারণত
\textenglish{\ttfamily None}
Return করে।

\begin{codeoutputbox}{Implicit \textenglish{\ttfamily None}}

\noindent
\begin{minipage}[t]{0.52\linewidth}
\codeheading

\begin{english}
\ttfamily
def greet():\\
\hspace*{4mm}print("Hello")\\[0.5mm]
x = greet()\\
print(x)
\end{english}

\end{minipage}
\hfill
\begin{minipage}[t]{0.44\linewidth}
\outputheading

\begin{english}
\ttfamily
Hello\\
None
\end{english}

\medskip

Function-এর ভিতরের
\textenglish{\ttfamily print()}
প্রথমে
\textenglish{\ttfamily Hello}
দেখিয়েছে।

Explicit
\textenglish{\ttfamily return}
না থাকায় Return Value হলো
\textenglish{\ttfamily None}।

\end{minipage}

\end{codeoutputbox}

%================================================
% Four Forms
%================================================

\booksubheading{Function-এর সাধারণ চারটি রূপ}

\begin{tabularx}
{\linewidth}
{
|>{\bfseries\centering\arraybackslash}p{0.22\linewidth}
|>{\centering\arraybackslash}p{0.21\linewidth}
|Y|
}

\hline
\rowcolor{softBlue}

\textbf{Parameter}
&
\textbf{Return}
&
\textbf{উদাহরণ}
\\
\hline

নেই
&
নেই
&
Message Print করা।
\\
\hline

আছে
&
নেই
&
Input নিয়ে Result Print করা।
\\
\hline

নেই
&
আছে
&
নির্দিষ্ট Value Return করা।
\\
\hline

আছে
&
আছে
&
Input নিয়ে Result Return করা।
\\
\hline

\end{tabularx}

%================================================
% Four examples compact
%================================================

\begin{codeoutputbox}{১. Parameter নেই, Return নেই}

\noindent
\begin{minipage}[t]{0.53\linewidth}
\codeheading

\begin{english}
\ttfamily
def show():\\
\hspace*{4mm}print("Hello")\\[0.5mm]
show()
\end{english}

\end{minipage}
\hfill
\begin{minipage}[t]{0.43\linewidth}
\outputheading

\begin{center}
\textenglish{\ttfamily Hello}
\end{center}

Function কোনো Input নেয়নি এবং কোনো Value Returnও করেনি।

\end{minipage}

\end{codeoutputbox}

\begin{codeoutputbox}{২. Parameter আছে, Return নেই}

\noindent
\begin{minipage}[t]{0.53\linewidth}
\codeheading

\begin{english}
\ttfamily
def add(a, b):\\
\hspace*{4mm}print(a + b)\\[0.5mm]
add(5, 3)
\end{english}

\end{minipage}
\hfill
\begin{minipage}[t]{0.43\linewidth}
\outputheading

\begin{center}
\textenglish{\ttfamily 8}
\end{center}

দুটি Argument নেওয়া হয়েছে,
কিন্তু Result শুধু Print করা হয়েছে।

\end{minipage}

\end{codeoutputbox}

\begin{codeoutputbox}{৩. Parameter নেই, Return আছে}

\noindent
\begin{minipage}[t]{0.53\linewidth}
\codeheading

\begin{english}
\ttfamily
def get\_number():\\
\hspace*{4mm}return 10\\[0.5mm]
x = get\_number()\\
print(x)
\end{english}

\end{minipage}
\hfill
\begin{minipage}[t]{0.43\linewidth}
\outputheading

\begin{center}
\textenglish{\ttfamily 10}
\end{center}

Function কোনো Parameter নেয়নি,
কিন্তু
\textenglish{\ttfamily 10}
Return করেছে।

\end{minipage}

\end{codeoutputbox}

\begin{codeoutputbox}{৪. Parameter আছে, Return আছে}

\noindent
\begin{minipage}[t]{0.53\linewidth}
\codeheading

\begin{english}
\ttfamily
def add(a, b):\\
\hspace*{4mm}return a + b\\[0.5mm]
x = add(5, 3)\\
print(x)
\end{english}

\end{minipage}
\hfill
\begin{minipage}[t]{0.43\linewidth}
\outputheading

\begin{center}
\textenglish{\ttfamily 8}
\end{center}

Function Input নিয়েছে এবং Result Return করেছে।

Practical Program-এ এটি খুব কার্যকর।

\end{minipage}

\end{codeoutputbox}

%================================================
% Multiple values
%================================================

\booksubheading{একাধিক Value Return}

Python Function একাধিক Value একসঙ্গে Return করতে পারে।

\begin{codeoutputbox}{একাধিক Result Return}

\noindent
\begin{minipage}[t]{0.55\linewidth}
\codeheading

\begin{english}
\ttfamily
def calculate(a, b):\\
\hspace*{4mm}return a + b, a - b\\[0.5mm]
x, y = calculate(10, 5)\\
print(x)\\
print(y)
\end{english}

\end{minipage}
\hfill
\begin{minipage}[t]{0.41\linewidth}
\outputheading

\begin{english}
\ttfamily
15\\
5
\end{english}

\medskip

একই Function থেকে Sum ও Difference ফেরত এসেছে।

Python একাধিক Return Value-কে সাধারণত Tuple আকারে প্রদান করে।

\end{minipage}

\end{codeoutputbox}

%================================================
% Function from function
%================================================

\booksubheading{একটি Function থেকে অন্য Function Call}

\begin{codeoutputbox}{Function-এর ভিতরে Function Call}

\noindent
\begin{minipage}[t]{0.55\linewidth}
\codeheading

\begin{english}
\ttfamily
def square(n):\\
\hspace*{4mm}return n * n\\[0.5mm]
def show():\\
\hspace*{4mm}print(square(5))\\[0.5mm]
show()
\end{english}

\end{minipage}
\hfill
\begin{minipage}[t]{0.41\linewidth}
\outputheading

\begin{center}
\textenglish{\ttfamily 25}
\end{center}

\textenglish{\ttfamily show()}
Function-এর ভিতর থেকে
\textenglish{\ttfamily square(5)}
Call করা হয়েছে।

এভাবে বড় Program ছোট ছোট Function-এ ভাগ করা যায়।

\end{minipage}

\end{codeoutputbox}

%================================================
% Built-ins
%================================================

\booksubheading{কিছু গুরুত্বপূর্ণ Built-in Function}

\begin{tabularx}
{\linewidth}
{
|>{\bfseries\RaggedRight\arraybackslash}p{0.31\linewidth}
|Y|
}

\hline
\rowcolor{softBlue}

\textbf{Function}
&
\textbf{কাজ}
\\
\hline

\textenglish{\ttfamily print()}
&
Output প্রদর্শন করে।
\\
\hline

\textenglish{\ttfamily input()}
&
User-এর কাছ থেকে Input গ্রহণ করে।
\\
\hline

\textenglish{\ttfamily len()}
&
Length বা Item Count নির্ণয় করে।
\\
\hline

\textenglish{\ttfamily type()}
&
Data Type নির্ণয় করে।
\\
\hline

\textenglish{\ttfamily int()}
&
Integer-এ Conversion করে।
\\
\hline

\textenglish{\ttfamily float()}
&
Float-এ Conversion করে।
\\
\hline

\textenglish{\ttfamily str()}
&
String-এ Conversion করে।
\\
\hline

\textenglish{\ttfamily range()}
&
Number Sequence তৈরি করে।
\\
\hline

\end{tabularx}

%================================================
% Naming rules
%================================================

\booksubheading{Function Name লেখার নিয়ম}

Function Name-ও Variable Name-এর সাধারণ নিয়ম অনুসরণ করে।

\begin{codeoutputbox}{সঠিক ও ভুল Function Name}

\noindent
\begin{minipage}[t]{0.48\linewidth}
\codeheading

\textbf{সঠিক—}

\begin{english}
\ttfamily
calculate\_sum()\\
student\_info()\\
area()
\end{english}

\end{minipage}
\hfill
\begin{minipage}[t]{0.48\linewidth}
\outputheading

\textbf{ভুল—}

\begin{english}
\ttfamily
2sum()\\
student-info()
\end{english}

\medskip

Name Digit দিয়ে শুরু করা যাবে না এবং
Hyphen ব্যবহার করা যাবে না।

\end{minipage}

\end{codeoutputbox}

Python Keyword Function Name হিসেবে ব্যবহার করা যাবে না।

%================================================
% Argument count
%================================================

\booksubheading{Function Call-এ Argument সংখ্যা}

\begin{codeoutputbox}{Argument সংখ্যা মিলতে হবে}

\noindent
\begin{minipage}[t]{0.52\linewidth}
\codeheading

\begin{english}
\ttfamily
def add(a, b):\\
\hspace*{4mm}return a + b\\[0.5mm]
print(add(10, 20))
\end{english}

\end{minipage}
\hfill
\begin{minipage}[t]{0.44\linewidth}
\outputheading

\begin{center}
\textenglish{\ttfamily 30}
\end{center}

এখানে দুটি Parameter আছে এবং দুটি Argument দেওয়া হয়েছে।

Argument কম বা বেশি হলে Error হতে পারে,
যদি Default বা বিশেষ Parameter ব্যবস্থা না থাকে।

\end{minipage}

\end{codeoutputbox}

%================================================
% Modular relation
%================================================

\booksubheading{Function ও Modular Programming}

Function ব্যবহার করে বড় Program-কে ছোট ছোট Logical অংশে ভাগ করা যায়।

\begin{center}

\textbf{Main Program}
\(\longrightarrow\)
\textbf{Input Function}
\(\longrightarrow\)
\textbf{Calculation Function}
\(\longrightarrow\)
\textbf{Output Function}

\end{center}

এতে প্রতিটি Function একটি নির্দিষ্ট কাজ করে এবং পুরো Program
আরও সংগঠিত হয়। Function তাই \textbf{Modular Programming}-এর
একটি গুরুত্বপূর্ণ ভিত্তি।

%================================================
% Execution flow
%================================================

\booksubheading{Function-এর Execution Flow}

\begin{codeoutputbox}{Execution Flow}

\noindent
\begin{minipage}[t]{0.53\linewidth}
\codeheading

\begin{english}
\ttfamily
def greet():\\
\hspace*{4mm}print("Hello")\\[0.5mm]
print("Start")\\
greet()\\
print("End")
\end{english}

\end{minipage}
\hfill
\begin{minipage}[t]{0.43\linewidth}
\outputheading

\begin{english}
\ttfamily
Start\\
Hello\\
End
\end{english}

\medskip

Execution Flow—

\begin{center}

Main Program

\(\downarrow\)

Function Call

\(\downarrow\)

Function Body

\(\downarrow\)

Main Program

\end{center}

\end{minipage}

\end{codeoutputbox}

%================================================
% Common Mistakes
%================================================

\chFiveInput{topics/44_python_functions/python_extensions.tex}

\booksubheading{সাধারণ ভুল}

\begin{rememberbox}

\begin{itemize}

\item
\textenglish{\ttfamily def}
Keyword লিখতে ভুল করা।

\item Function Header-এর শেষে
\textenglish{\ttfamily :}
না দেওয়া।

\item Function Body সঠিকভাবে Indent না করা।

\item Function Define করলেই Body Execute হবে মনে করা।

\item Function Call করার সময়
\textenglish{\ttfamily ()}
না দেওয়া।

\item Parameter ও Argument গুলিয়ে ফেলা।

\item প্রয়োজনীয় Argument-এর সংখ্যা না মেলানো।

\item
\textenglish{\ttfamily print()}
ও
\textenglish{\ttfamily return}
-কে একই মনে করা।

\item
\textenglish{\ttfamily return}
-এর পরের Statement Execute হবে মনে করা।

\item Invalid Function Name ব্যবহার করা।

\end{itemize}

\end{rememberbox}

%================================================
% Final Reminder
%================================================

\begin{rememberbox}

\textbf{\color{deepOrange}মনে রাখবে}

\begin{tabularx}
{\linewidth}
{
>{\bfseries\RaggedRight\arraybackslash}p{0.25\linewidth}
Y
>{\bfseries\RaggedRight\arraybackslash}p{0.25\linewidth}
Y
}

Function
&
Reusable Code Block
&
\textenglish{\ttfamily def}
&
Function তৈরি
\\[0.5mm]

Function Call
&
\textenglish{\ttfamily function\_name()}
&
Parameter
&
Definition-এর Variable
\\[0.5mm]

Argument
&
Call-এর Actual Value
&
\textenglish{\ttfamily return}
&
Result ফেরত দেয়
\\[0.5mm]

\textenglish{\ttfamily print()}
&
Output দেখায়
&
No explicit return
&
\textenglish{\ttfamily None}
\\[0.5mm]

Built-in
&
আগে থেকেই তৈরি
&
User-defined
&
Programmer তৈরি করে
\\[0.5mm]

Function
&
Code Reuse
&
Modularity
&
Programকে ছোট অংশে ভাগ
\\

\end{tabularx}

\medskip

\begin{center}

\textbf{Define}
\(\longrightarrow\)
\textenglish{\ttfamily def}

\qquad

\textbf{Run}
\(\longrightarrow\)
\textbf{Function Call}

\medskip

\textbf{Input to Function}
\(\longrightarrow\)
\textbf{Argument}

\qquad

\textbf{Receive inside Function}
\(\longrightarrow\)
\textbf{Parameter}

\medskip

\textbf{Result Back}
\(\longrightarrow\)
\textenglish{\ttfamily return}

\end{center}

\end{rememberbox}

\small

%================================================
% 5.45
%================================================

